% !TeX program = xelatex
\documentclass[UTF8,11pt,a4paper,fontset=fandol]{ctexart}

\usepackage[margin=23mm,headheight=15pt]{geometry}
\usepackage{amsmath,amssymb,amsthm,bm,mathtools}
\usepackage{booktabs,tabularx,array,longtable}
\usepackage{tikz}
\usetikzlibrary{arrows.meta,positioning,calc}
\usepackage{hyperref,xurl,fancyhdr,enumitem}
\hypersetup{hidelinks}
\setlength{\emergencystretch}{3em}
\setlength{\parskip}{0.35em}
\linespread{1.04}
\setlist{nosep,leftmargin=2em}
\numberwithin{equation}{section}

\newtheorem{theorem}{定理}[section]
\newtheorem{proposition}[theorem]{命题}
\newtheorem{lemma}[theorem]{引理}
\theoremstyle{definition}
\newtheorem{definition}[theorem]{定义}
\newtheorem{assumption}[theorem]{假设}
\newtheorem{remark}[theorem]{说明}
\renewcommand{\proofname}{证明}

\newcommand{\R}{\mathbb R}
\newcommand{\E}{\mathbb E}
\newcommand{\Id}{\mathrm I}
\newcommand{\dd}{\mathrm d}
\newcommand{\T}{^{\mathsf T}}
\newcommand{\norm}[1]{\left\lVert#1\right\rVert}
\newcommand{\abs}[1]{\left\lvert#1\right\rvert}
\DeclareMathOperator{\rank}{rank}
\DeclareMathOperator{\Cov}{Cov}
\DeclareMathOperator{\diag}{diag}
\DeclareMathOperator{\Log}{Log}
\DeclareMathOperator{\Retr}{Retr}
\DeclareMathOperator{\argmin}{arg\,min}
\DeclareMathOperator{\grad}{grad}
\DeclareMathOperator{\clip}{clip}
\newcolumntype{Y}{>{\raggedright\arraybackslash}X}

\pagestyle{fancy}
\fancyhf{}
\fancyhead[L]{\small Anatomical GPS V4}
\fancyhead[R]{\small 可行实现族的任务保持运输}
\fancyfoot[C]{\thepage}

\begin{document}

\title{
Anatomical GPS V4：\\
形变任务实现族的路径条件运输\\[2mm]
\large Path-Conditioned Transport of Anatomically Indexed
Feasible Realization Families
}
\author{NMI研究构造稿}
\date{2026年10月7日}
\maketitle

\begin{abstract}
本稿在V3的解剖任务表示、有限隐式实现单元和连续扫描走廊基础上，
将研究对象明确为：沿给定人体几何变化，连续运输一族具有相同成像
语义的机器人实现，而非仅重放或修复一条旧轨迹。
解剖地址定义不变任务身份，当前物理度量决定切向、平面及短轴关系，
实际运动学、声窗和配置约束定义可行实现集合。
本文构造物理代价一致的任务保持提升，以可微IRD选择允许的冗余运动，
并给出局部形变方向可行集及其对偶障碍。
在明确的正则性、速度选择和有限管域条件下，证明共同提升流能够运输
整个初始实现族，并保持扫描截面的连续性。
进一步说明终点相同不意味着运输结果相同，给出一个可精确计算的
超声视角例子、一个边界导致整族运输终止的例子，以及与V3有限单元
结合的验证流程。
连接、流、切锥、凸优化与认证延拓均具有成熟理论基础；本文不声称
这些数学工具首次出现。拟议研究增量在于由内部解剖任务生成并有限
验证可行实现族的共同运输。概率目标后验和真实执行安全分别说明，
不把模型条件结论等同于患者安全。
\end{abstract}

\noindent\textbf{关键词：}
解剖任务表示；任务保持运输；形变；约束流形；
连续扫描；可微可达性；路径依赖。

\begin{center}
\fbox{\begin{minipage}{0.94\linewidth}\small
\textbf{交付与验证状态。}
本稿包含解析推导及可以逐式检查的例子。
本轮文件执行环境不可用，未进行LaTeX编译、数值脚本运行或机器人实验。
文中数值仅由给出的解析式计算，不是患者或机器人结果。
不得将V3已有的实验、单元测试或代码存在性记作V4的已完成验证。
\end{minipage}}
\end{center}

\clearpage
\setcounter{tocdepth}{1}
\tableofcontents
\clearpage

% ============================================================
\section{V4究竟比V3增加什么？}

V3已经定义了从少量自由参数到接触位置、探头姿态和机器人逆解的
有限隐式提升。它也已经给出了有限模型差异下的根存在性与复用条件。
因此，重新把旧可行集合称作纤维，或把旧轨迹称作截面，不能算作
V4的实质贡献。

V4增加的对象是一个\emph{共同运输规则}。
它需要回答：从同一初始实现族出发，在同一人体变化路径下，
每个实现如何移动，哪些实现能够一直保持任务，何时由于实际约束
而无法继续，以及不同变化顺序为什么可能产生不同执行结果。

\begin{table}[htbp]\centering\small
\begin{tabularx}{\linewidth}{p{24mm}YY}
\toprule
问题 & V3 & V4\\
\midrule
基本对象 & 一个context中的有限可行实现单元
& 沿context路径演化的可行实现族\\
选择规则 & 走廊搜索及单元内优化
& 共同的任务保持提升与受约束冗余选择\\
有限保证 & 每个参数盒中的唯一实现
& 整个初始族在一个有限形变区间中的共同流\\
失败含义 & 未认证、排除或无连接
& 进一步区分任务提升障碍、速率障碍及运输域终止\\
历史信息 & 模型版本、证书、门点
& 还要保存运输规则、路径及当前分支\\
\bottomrule
\end{tabularx}
\caption{V4不删除V3，而是把V3单元用作共同运输的坐标与验证工具。}
\end{table}

本文保持UltraPoC为独立的物理接触执行模块。
不新增DiT、diffusion、人体生物力学规则或一般RMPflow控制架构。
主体采用微分运动学的几何提升和凸约束速度选择；空间GP仅用于目标
信念和相应任务容差。

% ============================================================
\section{最近理论基础及不可声称的创新}

iTaSC已经通过几何闭环统一特征、不确定参数、测量和速度控制
\cite{itasc}。冗余机器人的局部微分逆、重复性和Lie括号条件具有
直接先例\cite{shamir,tchon}。在正则、适当的投影下，将总空间视为
局部纤维化空间属于微分拓扑基础；Ehresmann型结论需要额外的
适当性、无边界或边界相容条件，不能仅由Jacobian满秩推出
全局运输\cite{lee,rabier}。

可生存性和控制障碍函数研究边界上的可行速度与前向不变性
\cite{ames}。自运动流形方法已经沿任务轨迹组织冗余解和连续路径
\cite{smm}。参数化Krawczyk方法已经认证随参数变化的解路径
\cite{duff}。

因此，以下主张不能单独算新颖：
\begin{quote}
``定义一个fiber bundle''；
``伪逆加零空间梯度''；
``人体变化后沿约束求导''；
``使用CBF或区间Newton''；
``局部都满足任务但闭合路径不回到初始构型''。
\end{quote}

本稿的研究对象更具体：
\begin{quote}
\textbf{由不变解剖地址和真实成像关系生成的实现族，
如何被一个共同规则沿有限人体变化运输，
以及哪些失败能够由实际任务与能力条件明确诊断。}
\end{quote}

这是一项候选研究构造，不是已经完成的首创性证明。
若最终实现只是逐点逆运动学加常规安全过滤，不能兑现上述区别。

% ============================================================
\section{解剖地址与物理任务}

\subsection{地址不变不等于物理度量不变}

设参考解剖域为$\mathcal A$，局部地址为$\xi$，
当前形态或模型context为$c$。当前世界位置为
\begin{equation}
 p_c=\Psi_c(\xi),\qquad
 J_c=D_\xi\Psi_c,\qquad
 G_c=J_c^\top J_c.
 \label{eq:anatomy}
\end{equation}
同一模型材料位置可以保留地址$\xi$，但长度、角度及正交关系由
$G_c$决定。跨患者仅能讨论解剖语义对应，不能将其称作同一物质粒子，
更不能把旧患者的后验当作新患者真值。

给定解剖切向$v$，世界切向按$J_cv$运输；
协向量按$J_c^{-\top}\alpha$运输。
如果任务要求切面在当前物理空间中垂直于血管，
应取
\begin{equation}
 \boxed{\alpha_c=G_cv,\qquad
 J_c^{-\top}\alpha_c=J_cv.}
 \label{eq:typed}
\end{equation}
固定材料平面和固定短轴要求一般不是同一件事。

\subsection{不规定人体形变生成机制}

输入为当前估计对应及有效域，而不是手写的人体形变定律。
主理论条件于一条给定的模型路径$c(\tau)$。
它可以来自真实观测的形态序列，或明确声明的模型插值。

虚拟模型插值不等于真实患者经历的运动。
如果患者已经处于终态$c_1$，中间context仅用于计算最终扫描族，
不能将中间context对应的位姿当作当前患者上的有效执行参考。

分片线性FEM图、角度接缝、低阶曲线和姿态对数分支必须分域处理。
主定理只在所需导数存在且有界的规则区域使用。
没有映射或误差界时，可以生成候选，不能宣称得到真实患者保证。

% ============================================================
\section{由内部目标生成可行实现族}

\subsection{一个明确的短轴扫描任务}

设当前内部中心线为$p_c(s)$，$s\in I=[s_0,s_1]$为材料进度，
单位切向为
\begin{equation}
 t_c(s)=\frac{\partial_s p_c(s)}
 {\norm{\partial_s p_c(s)}}.
\end{equation}
皮肤局部嵌入为$S_c(u)$，$u\in\R^2$，接触点为$o=S_c(u)$。

使用声学框架：$y_P$为切片法向，$z_P$为正深度，
$x_P$为图内横向。要求目标位于主声束且取得短轴切片时，
接触点满足
\begin{equation}
 \boxed{
 f(s,u,c)=t_c(s)^\top[S_c(u)-p_c(s)]=0.
 }
 \label{eq:shortaxis}
\end{equation}
令
\begin{equation}
 b=\frac{p_c(s)-S_c(u)}
 {\norm{p_c(s)-S_c(u)}}.
\end{equation}
在式\eqref{eq:shortaxis}上，$b\perp t_c$，因此
\begin{equation}
 R_P=[t_c\times b,\ t_c,\ b]\in SO(3).
\end{equation}

为在求根盒和导数计算中使用，应采用离流形的光滑正交延拓：
\begin{equation}
 \widetilde b=
 \frac{(\Id-t_ct_c^\top)[p_c-S_c]}
 {\norm{(\Id-t_ct_c^\top)[p_c-S_c]}}.
\end{equation}
要求分母在所用区域有正下界。不能将不正交矩阵直接送入
$SE(3)$对数。

\subsection{把机器人实现放进同一个约束}

令$q$包含机械臂关节与rail，$x=(u,q)$。
声学框架与FK工具框架之间使用实际标定。
联合任务方程为
\begin{equation}
 h(s,x,c)=
 \begin{bmatrix}
 f(s,u,c)\\
 \Log(T_P(s,u,c)^{-1}T_{\rm FK}(q))^\vee
 \end{bmatrix}=0.
 \label{eq:task}
\end{equation}
关节限位、碰撞、声窗、深度及模型有效范围写为
\begin{equation}
 g_j(s,x,c)\geq0.
\end{equation}
对七轴臂加rail，$x$有十个局部坐标，式\eqref{eq:task}
通常有七个独立等式，所以固定$s,c$时剩余三维自由度。
这只是规则区域的局部维数，不是全局拓扑保证。

\begin{definition}[解剖索引的可行实现族]
\begin{equation}
 \boxed{
 \mathcal F_c(s)=
 \{x:h(s,x,c)=0,\quad g_j(s,x,c)\geq0\}.
 }
 \label{eq:family}
\end{equation}
总空间及投影为
\begin{equation}
 \mathcal E=\{(s,c,x):x\in\mathcal F_c(s)\},
 \qquad
 \pi(s,c,x)=(s,c).
\end{equation}
\end{definition}

如果$D_xh$满行秩，无不等式边界的等式部分在局部具有规则纤维结构。
带边界、断开分支或维数改变的整个$\mathcal E$不自动是全局纤维丛。

\begin{definition}[扫描截面]
固定context$c$下，一条扫描是连续映射
$\sigma_c:I\to\mathcal X$，满足$\sigma_c(s)\in\mathcal F_c(s)$。
\end{definition}

% ============================================================
\section{运输不是由两个终点唯一决定的}

指定context路径$\gamma:\tau\mapsto c(\tau)$，并选定一个速度规则
$V(s,x,\tau)$。对固定解剖地址$s$，定义
\begin{equation}
 \partial_\tau x=V(s,x,\tau),\qquad x(0)=x_0.
 \label{eq:flow}
\end{equation}
若全过程位于相应可行族中，则得到
\begin{equation}
 \boxed{
 \Pi_{\gamma,s}^{\tau}(x_0)=x(\tau).
 }
 \label{eq:transport}
\end{equation}

必须在记号中保留$\gamma$及选择规则。
仅给出$c_0,c_1$通常不能决定唯一运输。
同一可行族也允许许多不同的零空间运动。
某条规则的运输失败，不等于所有任务保持运输均不存在。

\begin{definition}[所选规则的运输域]
\begin{equation}
 \mathcal D_{\gamma,s}
 =
 \{x_0\in\mathcal F_{c(0)}(s):
 \Pi_{\gamma,s}^{\tau}(x_0)
 \text{在全部 }0\leq\tau\leq1\text{ 上存在且可行}\}.
\end{equation}
\end{definition}

本文构造其可认证内域。
没有全域覆盖和排除证明时，不声称算出了最大运输域。

\begin{figure}[htbp]\centering
\begin{tikzpicture}[
 >=Latex,
 box/.style={draw,rounded corners,align=center,
 text width=35mm,minimum height=15mm,font=\small}]
 \node[box] (a) at (0,0) {源任务实现族\\$\mathcal F_{c_0}(s)$};
 \node[box] (b) at (5,0) {共同任务保持提升\\$\partial_\tau x=V$};
 \node[box] (c) at (10,0) {运输后的实现子族\\$\Pi_{\gamma,s}(\mathcal D_{\gamma,s})$};
 \draw[->] (a)--(b);
 \draw[->] (b)--(c);
 \node[align=center,font=\small,text width=125mm]
 at (5,-1.6) {解剖地址$s$固定；任务语义固定；
 选择规则、路径与可行边界共同决定实际运动。};
\end{tikzpicture}
\caption{V4运输一族实现，而非为各个$s$独立选取终态IK解。}
\end{figure}

% ============================================================
\section{物理代价一致的几何提升}

\subsection{任务微分关系}

记基变量$\beta=(s,c)$、其变化方向为$w=\dot\beta$，
并定义
\begin{equation}
 A=D_xh,\qquad B=D_\beta h.
\end{equation}
任何任务保持速度$v=\dot x$满足
\begin{equation}
 \boxed{Av+Bw=0.}
 \label{eq:tangent}
\end{equation}

此式是微分运动学基础，不因加入解剖context而自动成为新理论。
真正需要确定的是：在这一仿射速度空间中，采用哪个共同选择。

\subsection{不能在随人体移动的坐标中盲目最小化坐标速度}

$x=(u,q)$中的$u$可能是身体附着坐标。此时$\dot u=0$不表示接触点
在世界中不动。为定义物理代价，使用
\begin{equation}
 \Xi(\beta,x)=(S_c(u),q),\qquad
 D=D_x\Xi,\qquad e_w=D_\beta\Xi\,w.
\end{equation}
选择固定的物理标度矩阵$W\succ0$，最小化
\begin{equation}
 \frac12\norm{Dv+e_w}_W^2.
\end{equation}
当$D$满列秩时，令
\begin{equation}
 G=D^\top WD\succ0,\qquad
 v_{\rm ref}=-G^{-1}D^\top We_w.
\end{equation}
上式等价于最小化$\frac12\norm{v-v_{\rm ref}}_G^2$加常数。

\begin{proposition}[基准任务保持提升]
设$A$满行秩。定义
\begin{equation}
 A_G^\dagger=G^{-1}A^\top(AG^{-1}A^\top)^{-1},
 \qquad N=\Id-A_G^\dagger A.
\end{equation}
则最小物理代价的任务保持速度为
\begin{equation}
 \boxed{
 v_h=v_{\rm ref}
 -A_G^\dagger(Av_{\rm ref}+Bw).
 }
 \label{eq:horizontal}
\end{equation}
所有任务保持速度为$v=v_h+N\zeta$。
\end{proposition}

\begin{proof}
用拉格朗日乘子最小化二次代价并施加式\eqref{eq:tangent}，
消去乘子得到式\eqref{eq:horizontal}。
$AN=0$，且$v_h-v_{\rm ref}$在$G$度量下与$\ker A$正交，
因此该解唯一并具有最小代价。
\end{proof}

在独立于context的物理坐标中，若$v_{\rm ref}=0$，
式\eqref{eq:horizontal}退化为$-A_G^\dagger Bw$。
一般情况下不能省略移动坐标产生的漂移。

\subsection{表示一致性}

对context相关换图$x'=\varphi(x,\beta)$，速度满足
\begin{equation}
 v'=D_x\varphi\,v+D_\beta\varphi\,w.
\end{equation}
同时运输物理映射$\Xi$、任务$h$及约束后，
上述优化描述同一个物理速度，因此结果一致。
只对矩阵使用逆转置、却遗漏$D_\beta\varphi\,w$，不是正确的动态换图。

在无约束偏好、$v_h$对$w$线性且$w=0$时$v_h=0$的规则区域，
这给出一个Ehresmann连接的具体选择。
加入速度限制、单向能力偏好或活动集后，所用规则通常是受约束提升，
不再是线性的几何连接；不能混用二者的结论。

% ============================================================
\section{IRD只选择允许的运动，不定义真实可行性}

设$\rho(\beta,x)$是冻结参数的IRD能力代理。
$d_x\rho$是协向量，其度量梯度为$G^{-1}d_x\rho$。

无不等式限制时，可以选
\begin{equation}
 v=v_h+\kappa N G^{-1}d_x\rho,\qquad \kappa\geq0.
 \label{eq:irdlift}
\end{equation}
由于$N$是$G$正交投影，
\begin{equation}
 d_x\rho(v-v_h)
 =
 \kappa\norm{NG^{-1}d_x\rho}_G^2\geq0.
\end{equation}

这只表示能力引导相对于同一基准运输具有非负的瞬时增量。
完整变化率
\begin{equation}
 \dot\rho=D_\beta\rho\,w+d_x\rho(v_h)
 +\kappa\norm{NG^{-1}d_x\rho}_G^2
\end{equation}
仍可能为负，因为人体变化可能迫使机器人进入能力更差的区域。

用户IRD代码将其明确定义为可达裕量代理，而非碰撞证书。
真实可行族由实际运动学、关节约束与适用的安全模型定义。
除非另有保守误差证明，$\rho$只能指导选择，不能单独证明可行。

若希望在$w=0$时运输保持恒等，可将能力增益设成
$\kappa=\kappa_0\norm{w}_{G_{\rm base}}$。
但同号能力上升在反向context路径中不自动反向，
因此反走人体路径一般不是原运输的逆。

% ============================================================
\section{哪些形变方向能够被当前实现跟随？}

\subsection{几何切锥与带速率的可行集合}

在活动边界$g_j=0$处，任务保持和不越界要求
\begin{equation}
 Av=-Bw,\qquad
 D_xg_j\,v+D_\beta g_j\,w\geq0.
 \label{eq:activecone}
\end{equation}
在等式和活动不等式满足适当约束资格条件时，
这是总可行集合切锥的线性表达。

\begin{definition}[局部允许的基变化方向]
\begin{equation}
 \mathcal D_{\rm geom}(\beta,x)
 =
 \{w:\exists v\text{ 满足式\eqref{eq:activecone}}\}.
\end{equation}
\end{definition}
没有速度上限时它是锥。
加入$\dot q$等有界速率后得到$\mathcal D_{\rm rate}$，
一般不再是锥，不能继续利用任意比例放大的性质。

这个对象区分：
\emph{当前目标有解}与\emph{当前实现能沿指定变化继续移动}。
后者是运输所需的条件。

\subsection{实现用的速度QP}

为在内部也控制接近边界的速率，使用
\begin{equation}
 D_xg_jv+D_\beta g_jw\geq-\alpha_j g_j,
 \qquad \alpha_j>0.
 \label{eq:barrier}
\end{equation}
再加入实际速度范围$Dv\leq d$，记所得凸集合为$\mathcal K(\beta,x,w)$。

\begin{equation}
 \boxed{
 \begin{aligned}
 v^\star=\argmin_{v\in\mathcal K}\quad&
 \frac12\norm{v-v_h}_G^2
 -\kappa\,d_x\rho(v-v_h).
 \end{aligned}
 }
 \label{eq:qp}
\end{equation}

$G\succ0$保证非空闭凸可行集上的唯一最优解。
但唯一不等于在参数上自动光滑。
整族流定理还需要局部Lipschitz选择；固定活动集的LICQ及适当正则性
是一种充分条件。活动集转换必须检查，不能因QP每一步求解成功就
宣称运输流全局可微。

也可以在有限规则覆盖上构造平滑可行速度场：
各局部速度在同一点满足相同的仿射任务关系和凸速度约束时，
用非负、和为一的光滑权重组合，仍满足这些关系。
这要求局部可行选择有可用余量，不是任意混合两个机器人位姿。

\subsection{IRD相对基准的受约束改进}

令$v_0$为式\eqref{eq:qp}中$\kappa=0$的解。
以$v_0$作为可行比较点，由最优性和强凸性得到
\begin{equation}
 \kappa\,d_x\rho(v^\star-v_0)
 \geq \frac12\norm{v^\star-v_0}_G^2\geq0.
\end{equation}
它给出相对于同一约束下最小改动规则的局部能力增量；
不证明IRD的全局最优，也不证明未来整条路径都改善。

\subsection{对偶障碍，而不是仅返回求解失败}

将局部条件写成
\begin{equation}
 Av=r(w),\qquad Qv\leq d(w).
\end{equation}
根据Farkas型替代关系，不可行当且仅当存在
$\lambda$及$\mu\geq0$使
\begin{equation}
 \boxed{
 A^\top\lambda+Q^\top\mu=0,\qquad
 \lambda^\top r(w)+\mu^\top d(w)<0.
 }
 \label{eq:dual}
\end{equation}
该见证可以由线性规划求出并独立检查。

\begin{proof}
若存在可行$v$，则
$0=\lambda^\top Av+\mu^\top Qv
\leq\lambda^\top r+\mu^\top d<0$，矛盾。
反向是线性可行系统的分离定理。
\end{proof}

必须区分：
对\emph{真实活动边界切条件}和速度界的见证，否定该状态下所要求的
可微提升；对式\eqref{eq:barrier}的见证，可能只否定这个保守控制选择。
两者都不自动证明终态纤维为空，也不排除提前改变冗余状态后成功。

% ============================================================
\section{核心有限结果：运输整个初始实现族}

\begin{assumption}[共同运输的规则管域]
给定$c:[0,1]\to\mathcal C$及紧致初始参数域$I\times K$。
初始族$x_0(s,\zeta)$满足
$x_0(s,\zeta)\in\mathcal F_{c(0)}(s)$。
在包含其演化的共同有限管域中：
\begin{enumerate}
\item $h$和$g_j$具有所需连续导数，$D_xh$满行秩且远离所选图的奇异；
\item 选择场$V(s,x,\tau)$局部Lipschitz于$x$，连续依赖$(s,\tau)$；
\item $h_xV+h_c c'=0$；
\item $g_{j,x}V+g_{j,c}c'\geq-\alpha_jg_j$；
\item 解不会离开已验证的紧管域；该条件由边界条件或第
\ref{sec:validated}节的有限包含检查建立，而不是从一次数值积分猜测。
\end{enumerate}
\end{assumption}

\begin{theorem}[有限整族运输]
在上述条件下，每个$(s,\zeta)\in I\times K$存在唯一解
\begin{equation}
 X(s,\zeta,\tau),\qquad
 \partial_\tau X=V(s,X,\tau),\quad
 X(s,\zeta,0)=x_0(s,\zeta),
\end{equation}
并且对整个参数域和$0\leq\tau\leq1$，
\begin{equation}
 \boxed{
 X(s,\zeta,\tau)\in\mathcal F_{c(\tau)}(s).
 }
 \label{eq:familytheorem}
\end{equation}
解共同连续依赖$(s,\zeta,\tau)$。
固定$s$时，流对初始实现是单射；在相应光滑条件下，它是到其像的
光滑嵌入，而不必覆盖终态全部可行实现。
\end{theorem}

\begin{proof}
局部存在唯一性由ODE理论给出；紧管域和有界性使解延续到整个区间。
沿解有
$\dd h/\dd\tau=h_xV+h_cc'=0$，故等式保持。
由比较原理，
$g_j(\tau)\geq e^{-\alpha_j\tau}g_j(0)\geq0$。
连续依赖定理给共同连续性。
若两个不同初态在某时刻合并，局部反向唯一性将推出初态相同，
矛盾。光滑情况下，初态导数满足线性变分方程，其基本解矩阵可逆，
因此局部嵌入性保持。
\end{proof}

\begin{remark}
这不是一个新的Ehresmann定理。
它是把流、约束保持和有限验证条件作用于解剖任务实现族的构造。
约束边界可能只允许单向运输，终态也可能产生新的分支，
因此不能声称$\mathcal F_{c_0}(s)$与$\mathcal F_{c_1}(s)$全局微分同胚。
\end{remark}

\subsection{整条扫描的连续性不靠逐点重新选解}

若取初始扫描$\sigma_0(s)=x_0(s,\zeta(s))$，
共同运输得到
\begin{equation}
 \sigma_\tau(s)=X(s,\zeta(s),\tau).
\end{equation}
初始$\zeta(s)$连续则新扫描连续。

在一个共同光滑坐标中，令$M=\partial_s\sigma_\tau$。
则
\begin{equation}
 \partial_\tau M
 =
 V_xM+V_s.
\end{equation}
若$\norm{V_x}\leq L_x$、$\norm{V_s}\leq L_s$，有
\begin{equation}
 \boxed{
 \norm{M(\tau)}
 \leq e^{L_x\tau}\norm{M(0)}
 +\frac{L_s}{L_x}(e^{L_x\tau}-1).
 }
 \label{eq:sectionbound}
\end{equation}
$L_x=0$时使用极限$L_s\tau$。
这给出扫描形状连续变化的定量界；
没有这些正则性条件，独立逐点QP可能产生不可接受的分支或导数变化。

% ============================================================
\section{为什么不能只保存两个context终点？}

\subsection{连接曲率与路径依赖}

在无活动不等式、无独立垂直漂移的规则区域，
设基方向$\beta^i$对应提升系数$L_i(x,\beta)$。
总空间水平向量场为
\begin{equation}
 X_i=\partial_{\beta^i}+L_i.
\end{equation}
其括号的垂直分量为
\begin{equation}
 \boxed{
 \Omega_{ij}
 =
 \partial_iL_j-\partial_jL_i
 +D_xL_j\,L_i-D_xL_i\,L_j.
 }
 \label{eq:curvature}
\end{equation}
它描述小闭合基路径可能留下的冗余位移。

若希望存在一个局部、与路径无关的解码器，
并且它同时遵守这些水平选择，则混合偏导必须相容，
因此需要$\Omega_{ij}=0$。
在适当常秩和可积条件下可局部构造这种解码；
即便局部曲率为零，全局单值性仍可能受到拓扑影响。

上述理论与冗余机器人重复性研究直接相关
\cite{shamir,tchon}，不是本文发现。
V4使用它来约束GPS应该记录的内容：若选择最小运动或能力优化的
运输规则，不应再默认结果只依赖人体终态。

\subsection{一个完全解析的成像几何例子}

取皮肤平面$S(a,b)=(a,b,0)$，源血管为$(s,0,-d)$。
考虑解析映射
\begin{equation}
 \Phi_{c,\ell}(x,y,z)
 =
 \left(
 \frac{x+\ell}{1+c^2},
 \frac{c(x+\ell)}{1+c^2}+y,
 z
 \right).
 \label{eq:analyticmap}
\end{equation}
其行列式为$1/(1+c^2)>0$。
这是用于检验推导的模型，不是人体生物力学规则。

当前中心线和切向为
\begin{equation}
 p_{c,\ell}(s)=
 \left(
 \frac{s+\ell}{1+c^2},
 \frac{c(s+\ell)}{1+c^2},-d
 \right),\qquad
 t=\frac{(1,c,0)}{\sqrt{1+c^2}}.
\end{equation}
式\eqref{eq:shortaxis}等价于
\begin{equation}
 \boxed{h=a+cb-s-\ell=0.}
 \label{eq:toy}
\end{equation}

在物理接触坐标$(a,b)$采用欧氏最小速度。
令$D=1+c^2$，则$\ell$和$c$方向的水平提升分别为
\begin{equation}
 L_\ell=\frac{(1,c)^\top}{D},
 \qquad
 L_c=-\frac{b(1,c)^\top}{D}.
\end{equation}
计算得到
\begin{equation}
 \boxed{
 \Omega_{\ell c}=\frac{(c,-1)^\top}{(1+c^2)^2},
 \qquad [1,c]\Omega_{\ell c}=0.
 }
\end{equation}
所以剩余变化确实在任务零空间中。

从$c=\ell=0$出发，执行context矩形：
$\ell$增加$\Delta\ell$，$c$增加$C$，
再依次将$\ell$和$c$恢复到零。
直接积分得到
\begin{equation}
 \boxed{
 \Delta b=-\frac{\Delta\ell\,C}{\sqrt{1+C^2}},
 \qquad \Delta a=0.
 }
 \label{eq:holonomyexact}
\end{equation}
人体模型最终完全相同，成像任务始终满足，但最小速度运输留下了
不同接触实现。

例如$\Delta\ell=0.04\,\mathrm m$、$C=0.5$时，
$\Delta b\approx-17.89\,\mathrm{mm}$。
这是解析几何数值，不是机器人测量。

\subsection{路径独立与局部最小运动不能无条件兼得}

在式\eqref{eq:toy}中保持$b$不变，并令
$a=s+\ell-cb$，确实定义了一个局部路径独立的任务实现。
但它一般不是欧氏最小速度提升。

因此，不是``不存在可迁移表示''，
而是不同选择具有不同性质：
固定自由标签可以获得局部重复性；
最小物理运动或IRD偏好可以降低当前代价，
却可能产生路径依赖。
GPS必须明确保存采用哪种规则，不能同时无条件声称所有优点。

% ============================================================
\section{一个有限可行运输及其真实终止边界}

继续使用$\ell=0$的式\eqref{eq:toy}，附加局部允许接触域
\begin{equation}
 a\geq a_{\min},\qquad b\geq b_{\min},\qquad c\geq0.
\end{equation}
对固定$s$，可行实现存在的必要且充分条件为
\begin{equation}
 \boxed{
 c\leq c_\star(s):=\frac{s-a_{\min}}{b_{\min}},
 }
 \label{eq:critical}
\end{equation}
假设$s\geq a_{\min}$且$b_{\min}>0$。
它由$a+cb=s$直接得到，并可取$b=b_{\min}$达到。

在边界$a=a_{\min},b=b_{\min}$处，
\begin{equation}
 \dot a\geq0,\qquad \dot b\geq0,\qquad
 \dot a+c\dot b=\dot s-b_{\min}\dot c.
\end{equation}
故允许基方向必须满足
\begin{equation}
 \boxed{\dot s-b_{\min}\dot c\geq0.}
 \label{eq:toycone}
\end{equation}
固定地址运输时$\dot s=0$，不能继续沿$\dot c>0$前进。
这是一项真正的局部边界障碍，不是QP数值失败。

\subsection{无约束最小运动会提前失败}

取
\begin{equation}
 s=0.18\,\mathrm m,\quad
 a_{\min}=0.10\,\mathrm m,\quad
 b_{\min}=0.04\,\mathrm m,\quad
 b_0=0.10\,\mathrm m.
\end{equation}
无约束水平运输满足
\begin{equation}
 b(c)=\frac{b_0}{\sqrt{1+c^2}},\qquad
 a(c)=s-cb(c).
\end{equation}
它在$c=4/3$到达$a_{\min}$；
但可行族一直存在到$c_\star=2$。

所以，\emph{选定的无约束连接失效}不等于
\emph{任务已经不可实现}。

\subsection{一个运输整个初始区间的构造}

令
\begin{equation}
 s\in[0.18,0.42]\,\mathrm m,\qquad
 b_0\in[0.08,0.12]\,\mathrm m.
\end{equation}
定义共同运输
\begin{equation}
 \boxed{
 \begin{aligned}
 b(c;b_0)&=0.04+(b_0-0.04)(1-c/2)^2,\\
 a(s,c;b_0)&=s-cb(c;b_0).
 \end{aligned}
 }
 \label{eq:construct}
\end{equation}
对$0\leq c<2$，任务等式精确成立。
又因为
\begin{equation}
 a-0.10
 \geq
 (1-c/2)
 \left[
 0.08-(b_0-0.04)c(1-c/2)
 \right]
\end{equation}
且$c(1-c/2)\leq1/2$，
对$c\leq1.8$有统一下界
\begin{equation}
 a-0.10\geq0.004\,\mathrm m,\qquad
 b-0.04\geq0.0004\,\mathrm m.
\end{equation}
因此，整个$s$区间和整个初始$b_0$区间都被共同运输，
不是只检查一个初态或若干采样点。

它由同一个场生成：
\begin{equation}
 \frac{\dd b}{\dd c}
 =-\frac{b-b_{\min}}{1-c/2},\qquad
 \frac{\dd a}{\dd c}
 =-b-c\frac{\dd b}{\dd c}.
\end{equation}
在$c\leq1.8$时该场规则且有界。

\subsection{为什么不能把保证延伸到临界点？}

由式\eqref{eq:construct}，
\begin{equation}
 \frac{\partial b(c;b_0)}{\partial b_0}
 =(1-c/2)^2.
\end{equation}
在$c=2$时整个初始$b_0$区间合并到$b_{\min}$，
因此不再是单射的正则运输。
相应速度场的状态导数也在$c\to2$时失去统一有界性。

注意任务Jacobian$[1,c]$一直满行秩。
失效来自\emph{不等式边界使可行族塌缩}，
不是等式Jacobian必须先变奇异。
对于最小$s$，$c>2$时该局部允许域内的实现确实为空。

本例没有建模真实机械臂、实际IRD网络或组织声学。
它是V4主张可被解析检验的有限实例。

% ============================================================
\section{怎样与V3有限单元结合实现}
\label{sec:validated}

\subsection{V3是运输坐标，而不是被删除的旧方案}

V3提供局部自由变量$z$、从属变量$y$和认证提升
$y=Y(z,c)$。
V4在这些单元中表达共同运输场，再沿context方向推进整个初始族。

不能在每个$s$独立重新选一个终态根，然后宣称得到了共同运输。
单元之间必须保留V3的物理门点和分支一致性。

\subsection{有限步的共同包含检查}

设某一步起点族位于$X_0$，时间区间为$[0,\Delta]$，
候选管盒为$X$。
对所有相关$s$及context参数，用包含函数得到
$V(X,[0,\Delta])$和Lipschitz上界$L$。
若
\begin{equation}
 \boxed{
 X_0+[0,\Delta]\,V(X,[0,\Delta])
 \subseteq X,\qquad
 \Delta L<1,
 }
 \label{eq:picard}
\end{equation}
则Picard算子在相应函数管域内为收缩，
为每个起点和参数提供唯一的有限步解及共同包含范围。

同时检查该管域上的任务提升关系、约束导数和模型有效性。
其中任务等式可通过V3的精确隐式解码保持；
数值ODE残差必须独立计入误差预算。
用网格最大值替代真正包含界不构成认证。

验证式积分和参数化根追踪有成熟基础\cite{duff}。
本稿不声称一般高维问题具有多项式复杂度。

\subsection{具有数值残差时}

若近似路径满足
\begin{equation}
 \norm{\dot{\widehat x}-V(s,\widehat x,\tau)}
 \leq\delta_V,
 \qquad \norm{\widehat x(0)-x(0)}\leq e_0,
\end{equation}
且$\norm{V_x}\leq L$，则
\begin{equation}
 \norm{\widehat x(\tau)-x(\tau)}
 \leq e^{L\tau}e_0+
 \frac{\delta_V}{L}(e^{L\tau}-1).
\end{equation}
实际配置不等式需按其Lipschitz界收紧相应余量。
该误差是计算误差，不是患者模型误差或真实跟踪误差；
三者必须分别处理。

\subsection{算法}

\begin{enumerate}
\item 读取Anatomical GPS中的目标、任务类型、源实现族及分支；
      用当前材料映射生成$h,g$，不重放固定世界位姿。
\item 指定context路径及其物理或虚拟含义；仅有两个终点时，
      明确采用的插值，不能假装它是真实人体运动。
\item 在V3单元中构造$A,B,G,v_{\rm ref}$及IRD梯度，
      求可行提升或可验证的平滑近似。
\item 通过式\eqref{eq:picard}运输整个起点族；
      收紧后仍满足约束则存储新单元和运输映射。
\item 若图或约束资格条件失效，换图或缩短步长；
      若存在真实活动切条件的对偶障碍，记录局部障碍。
\item 只有存在经过验证的连接器，才能转移到另一实现分支；
      终态另有IK解不意味着可以无缝跳过去。
\item 输出运输成功域、剩余未判定域、局部障碍及适用模型版本。
      全局不可行仅由独立完整排除证据发布。
\end{enumerate}

% ============================================================
\section{从模型运输到实际扫描}

\subsection{扫描时间和运输参数不能混同}

构造好的截面为$\sigma(s,\tau)$。
若实际机器人沿$s(t)$执行，同时context确实沿$\tau(t)$变化，
则
\begin{equation}
 \dot x
 =
 \sigma_s\dot s+\sigma_\tau\dot\tau.
\end{equation}
在相应二阶正则条件下，
\begin{align}
 \ddot x={}&
 \sigma_{ss}\dot s^2+
 2\sigma_{s\tau}\dot s\dot\tau+
 \sigma_{\tau\tau}\dot\tau^2\\
 &+\sigma_s\ddot s+\sigma_\tau\ddot\tau.
\end{align}
这些导数和界用于选择可实现的扫描时序。

虚拟重规划时，先计算终态$\sigma(s,1)$，再在当前人体模型上执行；
不能在终态患者上执行$\tau<1$的虚拟中间位姿。
人体真实运动是外生的，降低扫描$\dot s$不等于降低人体$\dot c$。

\subsection{恢复反馈与DMP的关系}

实际执行误差可通过相同任务函数加入
\begin{equation}
 Av=-Bw-Kh.
\end{equation}
在理想模型和可执行条件下，
$\dot h=-Kh$，提供反应式恢复。
存在建模或执行扰动时，应分析$\dot h=-Kh+d$的有界误差，
不以DMP名称代替稳定性条件。

因此，本稿的主体属于\emph{任务函数微分运动学＋受约束几何提升}，
不是新的DMP或RMPflow。
DMP、G-DMP或ProDMP可以提供偏好或比较执行器，但不负责定义
可行实现族、运输域或跨分支连续性。
UltraPoC继续负责其独立的接触执行职责。

% ============================================================
\section{目标概率与未来管道：保留，但不篡改解剖}

\subsection{目标信念}

沿参考血管保持法丛残差过程，例如
\begin{equation}
 (\mathcal D_s+\lambda)^3e=L_w\,w_s,
\end{equation}
对应Matérn-$5/2$核
\begin{equation}
 K(s,s')=
 S_e\left(1+\lambda r+\frac{\lambda^2r^2}{3}\right)e^{-\lambda r},
 \qquad r=|s-s'|.
\end{equation}
其状态空间等价属于成熟GP理论\cite{hartikainen}。
分割观测必须来自实际切片交点、实际采集位姿与时间戳。
没有新观测时，重新满足几何闭环不等于增加解剖证据。

若需要更高阶名义几何导数，Matérn-$5/2$样本本身不自动满足所有
高阶运输条件。可以使用足够光滑的名义几何及独立误差包络；
不能偷偷把随机样本视为任意光滑。

\subsection{名义等式与概率容差}

名义任务可以采用后验均值构造$h=0$。
真实目标误差通过任务容差或独立的可信集合处理，例如
\begin{equation}
 \mathcal F_c^\delta(s)=
 \left\{
 x:
 \mathbb P\bigl(
 \norm{h(s,x,c;\theta)}_W\leq\varepsilon
 \mid\mathcal D
 \bigr)\geq1-\delta,\quad
 g_{\rm robot}\geq0
 \right\}.
\end{equation}
这不是要求同一个动作对全部随机目标都满足精确等式。

只在\emph{固定任务模型且不确定集合真正嵌套}时，
可以由集合包含推出鲁棒可行实现族变大。
一般Bayesian更新改变均值，可信集未必嵌套；
不能宣称“方差变小必然释放更多自由度”。

\subsection{目标分布与执行分布}

正确关系为
\begin{equation}
 p_\pi(\theta,X\mid\mathcal D)
 =
 p(\theta\mid\mathcal D)\,
 p_\pi(X\mid\theta,\mathcal D).
\end{equation}
第二项对每个$\theta$归一化，IRD不反向删除难以实现的目标假设。

分别为每条真实目标样本求一条最优轨迹，只得到条件候选族。
要称为未来可执行分布，必须使用同一因果策略，
在每个时刻仅接收该时刻能够获得的模拟观测。
有限节点的置信椭球也不自动成为整个连续时域的概率保证。

% ============================================================
\section{理论审查、验证设计和贡献边界}

\subsection{本轮可以明确支持什么？}

本文给出了：
\begin{enumerate}
\item 以解剖地址和成像关系索引的可行实现族；
\item 处理移动坐标漂移的物理代价一致提升；
\item 受任务、边界和速度限制的允许context方向及对偶障碍；
\item 在声明的规则有限管域上的整族运输定理；
\item 路径依赖的解析例子，以及不等式使整族运输终止的有限例子；
\item 使用V3有限单元验证共同运输的实施路线。
\end{enumerate}

\subsection{本轮没有完成什么？}

没有实现完整七关节机器人上的验证式运输器，
没有运行新代码或真实扫描，没有证明任意人体变化下安全，
也没有证明一般问题全局最优或多项式复杂度。
没有确认“可行实现运输”这一抽象本身首次出现。

\subsection{哪些对照最有决定性？}

保持相同目标、分割、机器人和任务要求，比较：
\begin{enumerate}
\item 固定局部姿态运输；
\item 新旧context独立逐点IK；
\item V3静态走廊重建／证书复用；
\item V4无不等式的最小代价提升；
\item V4受约束共同提升，分别有无IRD偏好。
\end{enumerate}
主要指标为整段任务残差、约束最小余量、分支连续性、
运输成功域、UNKNOWN覆盖和计算成本。
同一context终点的不同路径还应检验式\eqref{eq:holonomyexact}所描述的
路径依赖；不能要求一个本来非可积的最小代价规则无条件重复。

长血管主实验可固定患者形态、建立源扫描族，再改变到另一静态形态
并重建终态实现；动态人体跟随另设实验，避免混淆虚拟运输和物理运动。

\subsection{最强审稿意见}

审稿人可能指出：
``这是微分逆、CBF和认证延拓在一个新的任务约束上的应用。''

这项批评不能靠名称消除。
V4必须证明整族表示和运输，相比独立重规划与通用参数化追踪，
确实具有不可替代的计算或保证收益。
例如：共同运输域的有效计算、可靠的局部障碍分类、
较少重建代价下保持整段任务关系，以及人体路径依赖的正确处理。

若只实现了$\dot x=-h_x^\dagger h_c\dot c+N\nabla\rho$，
则尚未兑现V4；这条公式本身具有非常直接的成熟先例。

% ============================================================
\section{结论}

Anatomical GPS不应被理解成一张固定姿态表。
它定义的是不变解剖身份与任务关系，而当前人体、声窗和机器人状态
共同决定一族实现。

V4的构造是：
\begin{equation}
 \boxed{
 \mathcal F_{c_0}(s)
 \supseteq\mathcal D_{\gamma,s}
 \xrightarrow{\ \Pi_{\gamma,s}\ }
 \mathcal F_{c_1}(s),
 }
\end{equation}
其中运输只承诺在已认证初始域、指定context路径和选定规则上成立。
它可以运输一条扫描，也可以运输一整族扫描实现。

最重要的边界是：
\begin{quote}
\textbf{同一解剖地址不意味着唯一实现；
终态可达不意味着沿当前变化可跟随；
终点相同不意味着运输结果相同；
局部提升存在不意味着整族可以无限延拓。}
\end{quote}

这些边界使GPS保存的对象、IRD的作用和V3有限验证工具具有一致含义。
本稿提供的是可继续实现和审查的条件性构造，而不是以新术语替代已有
微分几何，也不是对真实患者安全的无条件声明。

% ============================================================
\begin{thebibliography}{99}

\bibitem{itasc}
J. De Schutter, T. De Laet, J. Rutgeerts, W. Decré, R. Smits,
E. Aertbeliën, K. Claes, and H. Bruyninckx.
Constraint-Based Task Specification and Estimation for Sensor-Based
Robot Systems in the Presence of Geometric Uncertainty.
\emph{The International Journal of Robotics Research},
26(5):433--455, 2007.
\url{https://doi.org/10.1177/0278364907078091}.

\bibitem{shamir}
T. Shamir and Y. Yomdin.
Repeatability of Redundant Manipulators: Mathematical Solution of the Problem.
\emph{IEEE Transactions on Automatic Control},
33(11):1004--1009, 1988.
\url{https://doi.org/10.1109/9.14412}.

\bibitem{tchon}
K. Tchoń.
Repeatable, extended Jacobian inverse kinematics algorithm for mobile manipulators.
\emph{Systems \& Control Letters}, 55(2):87--93, 2006.
\url{https://doi.org/10.1016/j.sysconle.2005.05.004}.

\bibitem{lee}
J. M. Lee.
\emph{Introduction to Smooth Manifolds}.
Second edition, Springer.
\url{https://doi.org/10.1007/978-1-4419-9982-5}.

\bibitem{rabier}
P. J. Rabier.
Ehresmann fibrations and Palais--Smale conditions for morphisms of Finsler manifolds.
\emph{Annals of Mathematics}, 146(3):647--691, 1997.
\url{https://doi.org/10.2307/2952457}.

\bibitem{ames}
A. D. Ames, S. Coogan, M. Egerstedt, G. Notomista,
K. Sreenath, and P. Tabuada.
Control Barrier Functions: Theory and Applications.
2019.
\url{https://arxiv.org/abs/1903.11199}.

\bibitem{smm}
M. Fabregat-Jaén, A. Peidró, M. Colombo, P. Rocco, and Ó. Reinoso.
Topological and spatial analysis of self-motion manifolds for global
redundancy resolution in kinematically redundant robots.
\emph{Mechanism and Machine Theory}, 210:106020, 2025.
\url{https://doi.org/10.1016/j.mechmachtheory.2025.106020}.

\bibitem{duff}
T. Duff and K. Lee.
Certified homotopy tracking using the Krawczyk method.
ISSAC 2024.
\url{https://arxiv.org/abs/2402.07053}.

\bibitem{hartikainen}
J. Hartikainen and S. Särkkä.
Kalman filtering and smoothing solutions to temporal Gaussian process regression models.
\emph{IEEE International Workshop on Machine Learning for Signal Processing}, 2010.
\url{https://doi.org/10.1109/MLSP.2010.5589113}.

\bibitem{v3}
NMI.
Anatomical GPS V3: 从不变解剖地址到任务保持的有限可行扫描.
用户项目研究稿，2026年10月7日。
本文引用其任务关系、有限隐式单元与分支门点作为V4输入，
不将V3未实现的完整验证器视为已存在。

\bibitem{ird}
NingMeng0903.
RealUS\_playground, ird\_playground/README.md.
核查版本：
\texttt{cd6b0cea1ae1983e55c1c732d0429b1bfd32218a}.
\url{https://github.com/NingMeng0903/RealUS_playground/blob/cd6b0cea1ae1983e55c1c732d0429b1bfd32218a/ird_playground/README.md}.

\end{thebibliography}
\end{document}